\section{Annuli and currents for a restricted HNN extension}
\label{ann:section}

We now isolate the combination statement used in the Magnus induction.
The restrictions on intersections in its hypotheses will provide a uniform
bound on the ranks of path stabilizers.  In this section we establish that
failure of hyperbolicity produces a bounded, nonzero trajectory of free-group
currents.  The subsequent sections rule out that trajectory and thereby
complete the combination argument.

\begin{theorem}[Restricted combination theorem]
\label{ann:combination}
Let $H$ be a torsion-free word-hyperbolic group, let $A,B\leq H$ be
finitely generated free quasiconvex subgroups, and let $\phi\colon A\to B$
be an isomorphism.  Set
\[
 G=\langle H,t\mid tat^{-1}=\phi(a)\ (a\in A)\rangle,
 \qquad P=A,\qquad I=A\cap B,\qquad J=t^{-1}It\leq P.
\]
Assume the following two conditions.
\begin{enumerate}[label=\textup{(\roman*)}]
\item For $M,N\in\{A,B\}$ and $g\in H$, the intersection
      $M\cap gNg^{-1}$ is at most cyclic unless $g\in MN$.
\item The groups $I,J$ are finitely generated, and each $E\in\{I,J\}$
      satisfies
      \begin{equation}
      \sum_{SpE\in S\backslash P/E}
        \bar r(E\cap p^{-1}Sp)\leq\bar r(S)
      \qquad\text{for every finitely generated }S\leq P.
      \label{ann:inertia-assumption}
      \end{equation}
\end{enumerate}
If $G$ contains no subgroup isomorphic to $BS(m,n)$ for any nonzero
integers $m,n$, then $G$ is word-hyperbolic.
\end{theorem}

All the intersection restrictions in this statement are hypotheses.
In particular, the assertion concerns a restricted class of HNN
extensions.  We retain the notation of Theorem~\ref{ann:combination}
throughout the combination argument.  Write
\[
 i\colon I\longrightarrow P,\quad i(c)=c,
 \qquad
 j\colon I\longrightarrow P,\quad j(c)=t^{-1}ct,
\]
so that $j(I)=J$.  Fix finite generating sets for $H$ and a free basis
for $P$.  We write $|\cdot|_H$ and $|\cdot|_P$ for the associated based
word lengths.  If $P=1$, then $G=H*\mathbb Z$ is hyperbolic, so we may
assume $P\ne1$ in what follows.

\subsection{Frames and stabilizers of paths}

Let $\mathcal T$ be the Bass--Serre tree of the HNN extension.  The height
homomorphism $G\to\mathbb Z$ sends $t$ to $1$ and $H$ to $0$.  Orient
the canonical edge $e$ from $t^{-1}H$ to $H$, and orient every translate
equivariantly.  Its stabilizer is $P$.  A \emph{frame} for the edge $he$
is the element $h\in G$; changing $h$ to $hp$, with $p\in P$, changes
the coordinates on its stabilizer but not the edge.  A path traverses an
edge positively if it follows this orientation.

\begin{lemma}[Neighboring edges and frames]
\label{ann:frames}
A reduced path in $\mathcal T$ with noncyclic pointwise stabilizer has
constant traversal sign.  The positive successors of $e$ that can occur
in such a path have the form $pte$, with $pI\in P/I$, and the positive
predecessors have the form $pt^{-1}e$, with $pJ\in P/J$.  Their common
stabilizers with $e$ are, respectively,
\[
 pIp^{-1}\quad\text{and}\quad pJp^{-1}.
\]
For a positive successor the coordinate transfer is
$c\mapsto t^{-1}p^{-1}cpt$; for a positive predecessor it is
$c\mapsto tp^{-1}cpt^{-1}$, on the indicated domains.
\end{lemma}

\begin{proof}
Every positively oriented edge leaving $H$ is $ate$ for some $a\in H$;
the possible edges are parametrized by $aB\in H/B$.  Its common
stabilizer with $e$ is $A\cap aBa^{-1}$.  If this intersection is
noncyclic, assumption~\textup{(i)} gives $a\in AB$.  Write $a=pb$ with
$p\in P$ and $b\in B$.  Since $bt=t\phi^{-1}(b)$, the edge $ate$ is
$pte$.  Moreover,
\[
 pte=qte
 \quad\Longleftrightarrow\quad
 t^{-1}q^{-1}pt\in A
 \quad\Longleftrightarrow\quad q^{-1}p\in A\cap B=I.
\]
The common stabilizer is
$P\cap pBp^{-1}=pIp^{-1}$, and conjugating by $(pt)^{-1}$ gives
the stated transfer.

A positive predecessor ending at $t^{-1}H$ is $t^{-1}ae$, with
$aA\in H/A$.  After conjugating its pair stabilizer by $t$, we obtain
$B\cap aAa^{-1}$.  Noncyclicity gives $a\in BA$; writing $a=br$ and
$p=t^{-1}bt\in P$ puts the predecessor in the form $pt^{-1}e$.
Two such edges agree exactly when
\[
 q^{-1}p\in P\cap t^{-1}Pt
 =t^{-1}(B\cap A)t=J.
\]
Their common stabilizer with $e$ is $pJp^{-1}$, and conjugation by
$(pt^{-1})^{-1}$ is the claimed transfer.

At a reversal of traversal sign, the consecutive edge stabilizers, in
the common vertex coordinates, are $M$ and $aMa^{-1}$ for the same
$M\in\{A,B\}$.  The edges are distinct because the path is reduced, so
$a\notin M$.  Assumption~\textup{(i)} makes their intersection at most
cyclic.  A noncyclic stabilizer of the whole path cannot lie in that
intersection.  Thus there are no reversals.
\end{proof}

We call a consecutive positive pair \emph{aligned} when its bridge lies
in the double coset $AB$ in the common vertex coordinates.  Reversing the
traversal gives the corresponding definition for a consecutive negative
pair.  Lemma~\ref{ann:frames} says that every path with noncyclic
stabilizer is aligned after its traversal sign is fixed.

\begin{lemma}[Rank bound for a fixed window]
\label{ann:window-bound}
Fix integers $\ell,r\geq0$.  Consider the positively traversed reduced
paths with edges indexed by $-\ell,\ldots,r$, with edge $0$ equal to $e$,
and with noncyclic pointwise stabilizer.  Then
\begin{equation}
 \sum_{[W]\in P\backslash\{\text{such paths}\}}
       \bar r\bigl(\operatorname{Stab}_G(W)\bigr)
 \leq\bar r(P).
 \label{ann:window-sum}
\end{equation}
Every stabilizer in this sum is finitely generated and free.  In
particular, there are at most $\bar r(P)$ such $P$-orbits.  Equivalently,
the same bound holds for $G$-orbits if the distinguished edge is allowed
to vary.
\end{lemma}

\begin{proof}
Build the fixed window by starting with $e$ and adding one edge at an
endpoint at a time.  The initial stabilizer is $P$.  Suppose $W$ is an
existing representative, $S=\operatorname{Stab}_G(W)\leq P$, and $h$
is a frame at the endpoint edge being extended.  In that frame put
$K=h^{-1}Sh\leq P$.

For a forward extension whose stabilizer is noncyclic,
Lemma~\ref{ann:frames} restricts the next edge to $hpte$, with
$pI\in P/I$.  The action of $S$ on these choices becomes left
multiplication by $K$.  Thus the extension orbits over the orbit of $W$
are indexed by the double cosets $KpI$, and the extended stabilizer is
isomorphic to $K\cap pIp^{-1}$.  Assumption~\textup{(ii)} gives
\[
 \sum_{KpI}\bar r(K\cap pIp^{-1})\leq\bar r(K)=\bar r(S).
\]
For backward extension use $J$ and its separately assumed inequality.
The stabilizers remain finitely generated: in the current edge
coordinates an extension takes the intersection of two finitely generated
subgroups of the finite-rank free group $P$.

The forgetful map from extended windows to the preceding windows is
$P$-equivariant.  Its fiber over the orbit of $W$ is precisely the set of
extensions modulo $S$, which is the double-coset set just considered.
Summing the displayed inequality over the existing representatives
therefore proves the induction step.  Windows with cyclic stabilizer
may be discarded, since extension only decreases a stabilizer.  This
proves~\eqref{ann:window-sum}.  Each noncyclic term contributes at least
one, proving the orbit bound.  Finally,
$\operatorname{Stab}_G(e)=P$, which identifies the two stated orbit sets.
\end{proof}

The shape $(\ell,r)$ is fixed in this lemma; no bound is asserted for the
union of windows of all lengths.  We shall also use the following direct
consequence of the same hypotheses.

\begin{lemma}
\label{ann:distinct-paths}
Two distinct aligned paths of the same fixed window shape with
distinguished edge $e$ and the same traversal sign have at most cyclic
common pointwise stabilizer.
\end{lemma}

\begin{proof}
At their first divergence from the common part, the common stabilizer,
in the shared endpoint frame, lies in
\[
 pEp^{-1}\cap qEq^{-1},\qquad pE\ne qE,
 \qquad E\in\{I,J\}.
\]
This intersection is at most cyclic.  Indeed, apply
\eqref{ann:inertia-assumption} with $S=E$.  The diagonal double coset
already contributes $\bar r(E)$, so every other term has reduced rank
zero.  Since $p^{-1}q\notin E$, the displayed intersection is one of
these off-diagonal intersections up to conjugacy.  If $E$ is cyclic or
trivial, the conclusion is immediate as well.
\end{proof}

\subsection{The algebraic annular-flaring criterion}

We state the external combination theorem with its quantifiers before
passing to limits.  This is important because the width bound in a
sequence of counterexamples will depend on the number of rows.

For a traversal sign $\epsilon\in\{+,-\}$ define the arriving and
departing maps into $H$ by
\[
 f_+(c)=c,\qquad f_-(c)=tct^{-1},\qquad
 l_+(c)=tct^{-1},\qquad l_-(c)=c.
\]
An \emph{algebraic annulus with $2m+1$ rows} consists of signs
$\epsilon_s\in\{+,-\}$ and elements $c_s\in P\setminus\{1\}$ for
$-m\leq s\leq m$, together with bridges $a_s\in H$ for
$-m\leq s<m$, satisfying
\begin{equation}
 f_{\epsilon_s}(c_s)
   =a_s l_{\epsilon_{s+1}}(c_{s+1})a_s^{-1}.
 \label{ann:row-equation}
\end{equation}
Its width is $\max_s|a_s|_H$ and its girth is $|c_0|_P$.  It is
\emph{essential} if a $+,-$ transition has $a_s\notin A$ and a $-,+$
transition has $a_s\notin B$.  There is no restriction at equal-sign
transitions.  It is $\lambda$-\emph{hyperbolic} if
\[
 \max\{|c_{-m}|_P,|c_m|_P\}\geq\lambda|c_0|_P.
\]
All these lengths are based word lengths.

For completeness, these definitions describe an actual reduced tree
path.  Write $f_{\epsilon_s}(c)=u^{-1}cu$ and
$l_{\epsilon_{s+1}}(c)=v^{-1}cv$, where $u,v\in\{1,t^{-1}\}$.
Starting with a frame $h_s$, choose
\begin{equation}
 h_{s+1}=h_su a_sv^{-1}.
 \label{ann:frame-equation}
\end{equation}
The two edges meet at $h_suH=h_{s+1}vH$, and
\eqref{ann:row-equation} gives
$h_sc_sh_s^{-1}=h_{s+1}c_{s+1}h_{s+1}^{-1}$.  Thus a single nontrivial
element fixes the developed path.  The four possible frame multipliers
are
\[
\begin{array}{c|cccc}
 (\epsilon_s,\epsilon_{s+1})&(+,+)&(+,-)&(-,+)&(-,-)\\ \hline
 u a_sv^{-1}&a_st&a_s&t^{-1}a_st&t^{-1}a_s.
\end{array}
\]
At the two reversing transitions the consecutive geometric edges agree
exactly when $a_s\in A$ or $a_s\in B$, respectively.  Essentiality is
therefore exactly absence of backtracking.

\begin{theorem}[Algebraic annular flaring]
\label{ann:external-flaring}
Let $H$ be word-hyperbolic, let $A,B\leq H$ be finitely generated,
and let $\phi:A\to B$ be an isomorphism such that both attaching
maps $A\to H$ are quasi-isometric embeddings. Set
$G=\langle H,t\mid tat^{-1}=\phi(a),\ a\in A\rangle$.
Then $G$ is word-hyperbolic if there exist $\lambda>1$ and an integer
$m\geq1$ such that, for every $\rho>0$, there is $h(\rho)$ for which
every essential algebraic annulus as defined above, with $2m+1$ rows,
width at most $\rho$, and girth at least $h(\rho)$, is
$\lambda$-hyperbolic.
\end{theorem}

This is the algebraic formulation of the Bestvina--Feighn combination
theorem in \cite[Theorem~1.2]{BestvinaFeighn1996Addendum}; we use its explicit
annular statement in \cite[Definitions~4.1--4.4 and Theorem~4.5]{Kapovich1999NonQuasiconvexity}.
The latter numbers refer to the arXiv version of that paper.  Our
$2m+1$ counts the successive edge groups, including both ends.  The
criterion applies here because $A$ and $B$ are quasiconvex in the
hyperbolic group $H$.

\begin{lemma}[Counterannuli at a fixed length]
\label{ann:counterannuli}
If $G$ is not hyperbolic, then for every $m\geq1$ there is a sequence
of essential annuli with $2m+1$ rows whose signs and bridges are fixed
throughout the sequence, such that
\[
 |c_0|_P\longrightarrow\infty,
 \qquad |c_{-m}|_P,|c_m|_P<2|c_0|_P.
\]
For each fixed $s\in[-m,m]$, the ratios $|c_s|_P/|c_0|_P$ are bounded.
The bounds and the fixed bridge tuple may depend on $m$.
\end{lemma}

\begin{proof}
Negate Theorem~\ref{ann:external-flaring} with $\lambda=2$.  For each
$m$ there is a finite width bound $\rho_m$ and counterannuli of
arbitrarily large girth with width at most $\rho_m$.  Choose such a
sequence with girth tending to infinity.  There are finitely many sign
tuples and finitely many bridge tuples of this fixed length and width,
so a subsequence fixes both tuples.

The maps $f_\epsilon,l_\epsilon$ are quasi-isometric embeddings.  The
row equation, with the bridges now fixed, gives constants $L\geq1$ and
$D\geq0$ such that
\[
 |c_{s+1}|_P\leq L|c_s|_P+D,
 \qquad |c_s|_P\leq L|c_{s+1}|_P+D
\]
at every row of this fixed sequence.  Iterating from row $0$ bounds all
the indicated ratios, since $|c_0|_P\to\infty$.
\end{proof}

\subsection{Currents and induction along an injection}

Let $F$ be a nontrivial finite-rank free group with a fixed free basis,
and let $T_F$ be its Cayley tree.  The space
\[
 \partial^2F=(\partial F\times\partial F)\setminus\{(\xi,\xi):
 \xi\in\partial F\}
\]
parametrizes oriented geodesic lines: $(\xi,\zeta)$ is directed from
$\xi$ to $\zeta$.  For a finite nondegenerate oriented segment $\gamma$
in $T_F$, let $\operatorname{Cyl}_F(\gamma)$ be the set of lines
containing it with that orientation.  These compact-open cylinders
form a basis for the topology.  Let $\mathcal L_F(v)$ be the set of
oriented lines passing through the vertex $v$.  It is a finite disjoint
union of the cylinders on edges directed out of $v$, hence is compact
open.

A \emph{current} on $F$ is a nonnegative, $F$-invariant, flip-invariant
Radon measure on $\partial^2F$.  Denote its cone by
$\operatorname{Curr}(F)$ and give it the topology of convergence on
compact-open cylinders.  Its norm is
\begin{equation}
 \|\mu\|_F=\frac12\mu\bigl(\mathcal L_F(1)\bigr).
 \label{ann:current-norm}
\end{equation}
This norm is additive on the cone.  It vanishes only on the zero
current, since translates of $\mathcal L_F(1)$ cover $\partial^2F$.
We omit the subscript when the group and basis are clear.

For $w\in F\setminus\{1\}$, the \emph{counting current} $\eta_w$
is the sum of point masses at the two orientations of the translates
of the axis of $w$, indexed by $F/\langle w\rangle$.  Indexing by
$\langle w\rangle$ retains the multiplicity when $w$ is a proper
power.  Directly counting the translated axes through $1$ gives
\begin{equation}
 \|\eta_w\|_F=|w|_{F,\mathrm{cyc}},
 \label{ann:counting-norm}
\end{equation}
where the right side is cyclic word length.  More generally, the
coordinate on a segment labeled by a reduced word $u$ counts the
occurrences of $u$ and $u^{-1}$ in the cyclic word for $w$, with
multiplicity.  For a cyclic word $z$, occurrences are tested at all
$|z|$ starting positions in its periodic repetition, including when
$|u|>|z|$.

\begin{lemma}[Compactness]
\label{ann:current-compactness}
The norm is continuous on $\operatorname{Curr}(F)$, and each closed
norm-bounded subset is compact.  In particular, every bounded sequence
of currents has a convergent subsequence.
\end{lemma}

\begin{proof}
Continuity follows because $\mathcal L_F(1)$ is compact open.  Every
compact subset of $\partial^2F$ is covered by finitely many translates
of this set.  A uniform norm bound therefore bounds the measures of
all compact subsets.  There are countably many segment cylinders.
Choose a diagonal subsequence on their measures.  The cylinder
identities are preserved: a cylinder is the finite disjoint union of
its extensions by one edge at either end.  On each compact-open
$\mathcal L_F(v)$ these identities define a finite measure, by
extension from its clopen cylinders, and the measures agree on
overlaps.  They consequently define a Radon measure on line space.
Invariance and flip invariance pass to the limit.  The continuous
norm retains its bound.  This proves sequential compactness; the
countable cylinder coordinates metrize the topology on a bounded
set, so compactness follows.
\end{proof}

We use the current map associated to an injection with its full
induction over cosets.  An ordinary boundary push alone would generally
fail to be invariant under the target group.  The construction below
agrees with \cite[Proposition--Definition~12.1]{KapovichCurrents2006}
in the nonabelian case, and also covers rank one.

\begin{lemma}[Induction of currents]
\label{ann:current-induction}
Let $\theta\colon Q\to F$ be an injection of nontrivial finite-rank
free groups.  It defines a continuous linear map
\[
 \theta_*\colon\operatorname{Curr}(Q)\longrightarrow
                    \operatorname{Curr}(F)
\]
which sends $\eta_w$ to $\eta_{\theta(w)}$ and sends every nonzero
current to a nonzero current.  For the fixed bases there are constants
$0<c_\theta\leq C_\theta<\infty$ such that
\begin{equation}
 c_\theta\|\nu\|_Q\leq\|\theta_*\nu\|_F
                     \leq C_\theta\|\nu\|_Q
 \qquad(\nu\in\operatorname{Curr}(Q)).
 \label{ann:induction-norm}
\end{equation}
Postcomposing $\theta$ with an inner automorphism of $F$ does not
change $\theta_*$.  Induction is functorial: if $\psi\colon F\to F'$
is another injection of nontrivial finite-rank free groups, then
\[
 (\psi\circ\theta)_*=\psi_*\circ\theta_*.
\]
\end{lemma}

\begin{proof}
Put $E=\theta(Q)$.  The subgroup $E$ is finitely generated and hence
quasiconvex in $F$.  Thus $\theta$ is a quasi-isometric embedding and
extends to a homeomorphism $\partial Q\to\Lambda E\subset\partial F$.
Push a current $\nu$ by this boundary map to an $E$-invariant measure
$\widehat\nu$ on $\partial^2\Lambda E$, viewed as a measure on
$\partial^2F$, and define
\begin{equation}
 \theta_*\nu=\sum_{pE\in F/E}p_*\widehat\nu.
 \label{ann:coset-induction}
\end{equation}
The summand depends only on the coset.  It remains to prove that this
sum is locally finite and continuous.

Let $T_E$ be the convex hull of $\Lambda E$ in $T_F$.  Its quotient
by $E$ is the finite unbased core graph of $E$.  For a fixed vertex
$v$, there are only finitely many cosets $pE$ for which $v\in pT_E$:
indeed such cosets correspond to $E$-orbits of vertices $p^{-1}v$ in
$T_E$.  A coset term in~\eqref{ann:coset-induction} that meets a
segment cylinder has its hull containing that segment.  Hence only
finitely many terms contribute to a fixed cylinder.  The preimage of
that cylinder in $\partial^2Q$ is compact open, because
$\partial^2\Lambda E$ is closed in $\partial^2F$ and the boundary
map is a homeomorphism onto it.  It follows that the sum is Radon
and that every cylinder coordinate depends continuously and linearly
on $\nu$.  The sum is $F$-invariant and flip-invariant by construction.

The coset decompositions in~\eqref{ann:coset-induction} combine the
index sets $F/E$ and $E/\theta(\langle w\rangle)$ into
$F/\langle\theta(w)\rangle$.  This proves the assertion on counting
currents, including proper powers.  The same coset calculation proves
functoriality.  Indeed, in $\psi_*\theta_*\nu$ the outer representatives
range over $F'/\psi(F)$ and the inner representatives over
$\psi(F)/\psi\theta(Q)$.  Their products form representatives for
$F'/\psi\theta(Q)$, and the boundary maps compose.  Thus the two
nonnegative sums give exactly $(\psi\circ\theta)_*\nu$.

A nonzero measure has nonzero
boundary push, and the summands are nonnegative, so its induction is
nonzero.  By Lemma~\ref{ann:current-compactness}, the norm-one set
in $\operatorname{Curr}(Q)$ is compact.  The continuous function
$\nu\mapsto\|\theta_*\nu\|_F$ on that set is positive, so its
minimum and maximum give~\eqref{ann:induction-norm}.  Finally,
conjugating $E$ and the boundary map simply reindexes the coset sum;
the resulting $F$-invariant current is unchanged.
\end{proof}

No injectivity of the linear map $\theta_*$ is asserted or needed.
Nonzero preservation and the two norm inequalities are the properties
that will be used.

\subsection{From based words to currents}

The girth of an annulus uses based length, whereas the norm of a
counting current uses cyclic length.  A conjugating tail can make
these lengths very different.  We bridge that difference by capping
the word in the intersection subgroup before applying either of its
two embeddings.

For reduced words $u,w$ over a fixed basis, let $N_u(w)$ count the
occurrences of $u$ as a contiguous subword of the linear word $w$.
Occurrences crossing an end are not counted.  Write
$N_u^{\mathrm{sym}}(w)=N_u(w)+N_{u^{-1}}(w)$.

\begin{lemma}[Capping linear words]
\label{ann:capping}
Fix an injection $\theta\colon Q\to F$ of nontrivial finite-rank
free groups with fixed bases.  For each nonempty reduced word $W$ in
$Q$, there is a suffix $b$ of length at most one such that $Wb$ is
cyclically reduced.  For every fixed nonempty reduced word $u$ in $F$
there is a constant $D_{\theta,u}$, independent of $W$, such that
\begin{equation}
 \left|\eta_{\theta(Wb)}
       \bigl(\operatorname{Cyl}_F([1,u])\bigr)
       -N_u^{\mathrm{sym}}\bigl(\operatorname{red}(\theta(W))\bigr)
 \right|\leq D_{\theta,u}.
 \label{ann:capping-frequency}
\end{equation}
There is also a constant $D_\theta$ independent of $W$ such that
\begin{equation}
 \bigl|\|\eta_{\theta(Wb)}\|_F-|\theta(W)|_F\bigr|
 \leq D_\theta.
 \label{ann:capping-length}
\end{equation}
Here $\operatorname{red}$ denotes the unique reduced representative.
The same cap $b$ may be used simultaneously for any finite collection
of fixed injections out of $Q$.
\end{lemma}

\begin{proof}
If $Q$ is cyclic, every nonempty reduced word is already cyclically
reduced, so take $b$ empty.  Otherwise there are at least four oriented
basis letters.  Choose a letter $b$ different from the inverse of the
last letter of $W$ and from the inverse of its first letter.  Then
$Wb$ is freely and cyclically reduced.  This choice depends only on
$Q$ and $W$, not on $\theta$.

The image in $T_F$ of a geodesic in $T_Q$, with the images of
successive vertices joined by geodesic segments, is a uniform
quasigeodesic.  Its constants depend only on $\theta$ and the bases.
The axis of $Wb$ in $T_Q$ passes through $1$.  Its image is therefore
a uniform quasigeodesic line passing through $1$, with endpoints the
fixed points of $\theta(Wb)$.  The Morse lemma gives a constant
$D$, independent of $W$, bounding the distance from $1$ to the
axis of $\theta(Wb)$ in $T_F$.  Accordingly, if
\[
 \operatorname{red}(\theta(Wb))=v z v^{-1}
\]
with $z$ cyclically reduced and the concatenation reduced, then
$|v|_F\leq D$.

Multiplication by $\theta(b)$ changes only a bounded suffix of
$\operatorname{red}(\theta(W))$: at most
$\max\{|\theta(x)|_F:x\text{ a basis letter of }Q\}$ letters can
be canceled, and at most that many can be added.  Removing $v$ and
$v^{-1}$ then deletes at most $2D$ additional letters.  For a fixed
word $u$, changing or deleting a bounded number of letters at the
ends changes its linear occurrence count by a bounded amount.
Passing from a linear word $z$ to its cyclic occurrence count changes
the count by at most $|u|-1$ at each chosen orientation, even when
$|z|<|u|$.  These observations prove
\eqref{ann:capping-frequency}, including the inverse occurrences.
The same bounded additions and deletions show that
$|z|_F=|\theta(W)|_F+O_\theta(1)$; now use
\eqref{ann:counting-norm} to obtain~\eqref{ann:capping-length}.
All constants can be chosen for each member of a finite collection
of injections while retaining the same cap.
\end{proof}

\begin{lemma}[A current limit for each finite window]
\label{ann:finite-current-window}
Suppose $G$ is not hyperbolic, and fix $m\geq1$.  There are currents
$\mu_s\in\operatorname{Curr}(P)$ for $-m\leq s\leq m$ with
\begin{equation}
 \|\mu_0\|=1,\qquad \|\mu_{-m}\|,\|\mu_m\|\leq2,
 \label{ann:finite-end-bounds}
\end{equation}
and a fixed essential sign and bridge tuple such that each consecutive
pair is induced by a current on its intersection subgroup.  More
precisely, for
\[
 K_s=f_{\epsilon_s}(P)\cap
             a_s l_{\epsilon_{s+1}}(P)a_s^{-1}
\]
there is $\nu_s\in\operatorname{Curr}(K_s)$ satisfying
\begin{equation}
 (\theta_{s,0})_*\nu_s=\mu_s,
 \qquad(\theta_{s,1})_*\nu_s=\mu_{s+1},
 \label{ann:finite-compatibility}
\end{equation}
where
\[
 \theta_{s,0}=f_{\epsilon_s}^{-1},\qquad
 \theta_{s,1}=l_{\epsilon_{s+1}}^{-1}
                \circ(c\mapsto a_s^{-1}ca_s).
\]
All the groups $K_s$ are nontrivial finitely generated free groups.
\end{lemma}

\begin{proof}
Take the sequence from Lemma~\ref{ann:counterannuli}, keeping $m$,
all signs, and all actual bridges fixed.  Write its elements as
$c_s^{(k)}$ and put $L_k=|c_0^{(k)}|_P\to\infty$.
Lemma~\ref{pre:quasiconvex-intersection} shows that $K_s$ is finitely
generated.  It is free because it lies in the free group
$f_{\epsilon_s}(P)$, and is nontrivial because it contains
$f_{\epsilon_s}(c_s^{(k)})$.  Fix a basis for each $K_s$.

Represent $f_{\epsilon_s}(c_s^{(k)})$ by a reduced word $W_s^{(k)}$
in the chosen basis of $K_s$.  Apply Lemma~\ref{ann:capping} to
choose a cap $b_s^{(k)}$, and form
\[
 \nu_s^{(k)}=L_k^{-1}\eta_{W_s^{(k)}b_s^{(k)}}.
\]
The injection $\theta_{s,0}$ is quasi-isometric between free groups.
Consequently $|W_s^{(k)}|_{K_s}$ is at most a fixed multiple of
$|c_s^{(k)}|_P$ plus a fixed constant.  The ratio bound from
Lemma~\ref{ann:counterannuli} shows that $\|\nu_s^{(k)}\|$ is
bounded.  There are only finitely many $s$, so after passing to a
subsequence all $\nu_s^{(k)}$ converge, to currents $\nu_s$.

Apply Lemma~\ref{ann:capping} to both injections from $K_s$.  The row
equation gives
\[
 \theta_{s,0}(W_s^{(k)})=c_s^{(k)},\qquad
 \theta_{s,1}(W_s^{(k)})=c_{s+1}^{(k)}.
\]
For every fixed word cylinder, the two induced normalized currents
therefore differ from the corresponding normalized symmetrized
linear row frequencies by $O(1)/L_k$.  At an interior row, the two
neighboring induced limits agree, since both are limits of the same
linear frequencies in $c_s^{(k)}$.  Denote the resulting current by
$\mu_s$; at the end rows use their single adjacent transitions.
Continuity from Lemma~\ref{ann:current-induction} gives
\eqref{ann:finite-compatibility}.

The length conclusion~\eqref{ann:capping-length} gives
\[
 \|\mu_s\|=\lim_k\frac{|c_s^{(k)}|_P}{L_k}
\]
along this subsequence.  The middle ratio is $1$ and the two end
ratios are less than $2$, proving~\eqref{ann:finite-end-bounds}.
The errors used here may depend on $m$, the actual bridges, the
bases of $K_s$, and the test cylinder.  They disappear because
$m$ and all those data were fixed before $L_k\to\infty$.
\end{proof}

\subsection{Finite transition types and a bilateral limit}

The preceding construction produces one finite current trajectory for
each $m$.  To pass to unbounded window lengths, we first replace the
actual bridges by a finite list of double-coset representatives.  This
replacement is made after the current limits, so it does not require
any bound on the original bridges that is uniform in $m$.

For each ordered sign pair $(\epsilon,\delta)$, choose representatives
$d$ for the double cosets
\[
 f_\epsilon(P)\backslash H/l_\delta(P)
\]
with nontrivial intersection
\[
 K_{\epsilon,\delta,d}
   =f_\epsilon(P)\cap d l_\delta(P)d^{-1}.
\]
There are finitely many, by Lemma~\ref{pre:double-coset-finiteness}
and torsion-freeness of $H$.  Use $d=1$ for the identity double coset
when its intersection is nontrivial.  For reversing sign pairs discard
that double coset, since it gives precisely the forbidden pinches.
Call the resulting finite set of triples
$\tau=(\epsilon,\delta,d)$ the \emph{transition types}, and write
\begin{equation}
 K_\tau=K_{\epsilon,\delta,d},\qquad
 \theta_{\tau,0}=f_\epsilon^{-1},\qquad
 \theta_{\tau,1}=l_\delta^{-1}\circ(c\mapsto d^{-1}cd).
 \label{ann:type-maps}
\end{equation}
Fix a free basis for every $K_\tau$.

\begin{lemma}[Uniformity of transition types]
\label{ann:finite-types}
Every consecutive pair of currents from
Lemma~\ref{ann:finite-current-window} can be expressed as
\[
 \mu_s=(\theta_{\tau_s,0})_*\nu_s,
 \qquad
 \mu_{s+1}=(\theta_{\tau_s,1})_*\nu_s
\]
for a current on one of the fixed groups $K_{\tau_s}$, with type
signs $(\epsilon_s,\epsilon_{s+1})$.  There are constants $C\geq2$
and $D\geq1$, independent of $m$, such that
\begin{equation}
 C^{-1}\|\mu_s\|\leq\|\mu_{s+1}\|\leq C\|\mu_s\|,
 \qquad \|\nu_s\|\leq D\|\mu_s\|.
 \label{ann:uniform-transition-bounds}
\end{equation}
If $K_\tau$ is noncyclic, the signs are equal and $d=1$.
For the positive aligned type the two maps are $i,j$, and for the
negative aligned type they are $j,i$.
\end{lemma}

\begin{proof}
Suppress the row index and write $f=f_\epsilon$, $l=l_\delta$.
Factor an actual bridge as $a=f(p)d l(q)$ with $p,q\in P$.
Conjugation by $f(p)$ identifies
$K_d=f(P)\cap d l(P)d^{-1}$ with
$K_a=f(P)\cap a l(P)a^{-1}$.  Under this identification, the maps
from $K_a$ to the two edge coordinates become
\[
 k\longmapsto p\theta_{\tau,0}(k)p^{-1},
 \qquad
 k\longmapsto q^{-1}\theta_{\tau,1}(k)q.
\]
By Lemma~\ref{ann:current-induction}, these inner automorphisms leave
the induced currents unchanged.  Transporting the
transition current along the group isomorphism therefore gives the
desired representation.  Essentiality is unchanged: for a reversal,
membership of the bridge in the common side subgroup depends only
on its double coset.

Each of the finitely many maps in~\eqref{ann:type-maps} has upper
and positive lower norm bounds from
Lemma~\ref{ann:current-induction}.  Comparing the two bounds for a
type and then taking maxima over the finite type set proves
\eqref{ann:uniform-transition-bounds}.

If $K_\tau$ is noncyclic, assumption~\textup{(i)} puts $d$ in the
identity double coset.  At a reversal that coset was discarded, so
the signs must be equal.  For equal signs the intersection is $I$,
and direct substitution in~\eqref{ann:type-maps} gives $i,j$ for
$(+,+,1)$ and $j,i$ for $(-,-,1)$.
\end{proof}

\begin{proposition}[The current trajectory forced by nonflaring]
\label{ann:trajectory}
Under the hypotheses of Theorem~\ref{ann:combination}, if $G$ is not
hyperbolic, there exist currents
$\mu_s\in\operatorname{Curr}(P)$ for $s\in\mathbb Z$, signs
$\epsilon_s\in\{+,-\}$, and transition types $\tau_s$ with signs
$(\epsilon_s,\epsilon_{s+1})$, such that
\begin{enumerate}[label=\textup{(\roman*)}]
\item $\|\mu_0\|=1$, every $\mu_s$ is nonzero, and
      $\sup_{s\in\mathbb Z}\|\mu_s\|<\infty$;
\item for every $s$ there is a current
      $\nu_s\in\operatorname{Curr}(K_{\tau_s})$ with
      \[
      \mu_s=(\theta_{\tau_s,0})_*\nu_s,
      \qquad
      \mu_{s+1}=(\theta_{\tau_s,1})_*\nu_s;
      \]
\item all types are essential: any reversing type uses a double
      coset outside the common side subgroup.
\end{enumerate}
The constants in~\eqref{ann:uniform-transition-bounds} apply to this
trajectory.  In particular, $\|\mu_s\|\geq C^{-|s|}$ for all $s$.
\end{proposition}

\begin{proof}
For each $m$ take one finite current trajectory from
Lemma~\ref{ann:finite-current-window}, express its transitions using
Lemma~\ref{ann:finite-types}, and put
\[
 D_m=\max_{-m\leq s\leq m}\|\mu_s\|.
\]
The constants $C,D$ in~\eqref{ann:uniform-transition-bounds} are
now independent of $m$.

Suppose first that $D_m$ is bounded on a subsequence $m\to\infty$.
Keep its original centers.  On each fixed finite window the row
currents have bounded norm, as do the transition currents by
\eqref{ann:uniform-transition-bounds}.  The sign and type sets are
finite.  Successively extract subsequences on the windows
$[-n,n]$, fixing the signs and types there and taking limits of the
finitely many row and transition currents.  A diagonal subsequence
gives currents and compatible signs and types for all $s\in\mathbb Z$.
Continuity of the induced maps preserves the two current equalities.
Norm continuity preserves $\|\mu_0\|=1$ and the uniform upper bound.

In the remaining case choose a subsequence with $D_m\to\infty$,
and choose an index $j_m$ where the maximum is attained.  Divide all
row and transition currents in that finite trajectory by $D_m$, and
reindex so that $j_m$ becomes $0$.  The new row norms are at most
$1$, and the central norm is $1$.  Both original end norms were
at most $2$.  Iterating the adjacent norm comparison therefore gives
\[
 D_m\leq2C^{j_m+m},\qquad D_m\leq2C^{m-j_m}.
\]
Thus
\begin{equation}
 \min\{j_m+m,m-j_m\}\geq\log_C(D_m/2)\longrightarrow\infty.
 \label{ann:recenter-distance}
\end{equation}
Every fixed bilateral window is eventually present in the recentered
trajectory.  The same diagonal extraction applies; the transition
norms remain uniformly bounded after normalization.  It produces
a bilateral trajectory with central norm $1$ and all row norms at
most $1$.

In either case, the finite-type equalities and their norm comparisons
survive the limit.  Applying the lower comparison repeatedly from
row $0$ gives $\|\mu_s\|\geq C^{-|s|}$, so none of the row currents
vanishes.  Compatibility of the two signs of a transition with its
neighboring rows survives because these discrete data were fixed
on each extracted finite window.  The forbidden reversing types
were excluded from the type set at the outset.  This proves all
three conclusions.
\end{proof}

There are two distinct limits in this proof.  At each fixed window
length, first send the original annular girth to infinity, with all
actual bridges fixed.  Only after those errors have disappeared do we
replace bridges by finitely many types and let the window length grow.
The quantitative estimate~\eqref{ann:recenter-distance} ensures that
recentering cannot lose either half of the trajectory, while continuity
of the norm preserves nonzero mass.  We have therefore reduced
Theorem~\ref{ann:combination} to excluding the bounded current
trajectory of Proposition~\ref{ann:trajectory} under BS-freeness.
